% SEGMENT OLP-0057-S01
% Part: propositional-logic
% Chapter: syntax-and-semantics
% Section: introduction

\documentclass[../../../include/open-logic-section]{subfiles}

\begin{document}

\olfileid{pl}{syn}{int}

\olsection{அறிமுகம்}

கூற்றுத் தருக்கவியலின் !!{formula}s, !!{propositional variable}s
மற்றும் $\lnot$, $\land$, $\lor$, $\lif$, $\liff$ ஆகிய கூற்று
இணைப்புகளைப் பயன்படுத்திக் கட்டமைக்கப்படுகின்றன. உள்ளுணர்வுப்படி,
!!a{propositional variable}~$p$ என்பது மெய்யோ பொய்யோ ஆன ஒரு
வாக்கியம் அல்லது கூற்றைக் குறிக்கிறது. ஒரு !!a{formula}வில் வரும்
!!{propositional variable}களின் ``மெய்மதிப்பு'' தீர்மானிக்கப்பட்டவுடன்,
அவற்றிலிருந்து கூற்று இணைப்புகளால் கட்டமைக்கப்பட்ட ஒவ்வொரு
!!a{formula}வின் மெய்மதிப்பும் தீர்மானிக்கப்படுகிறது. ஆகவே அதன்
பொருண்மை மெய்மதிப்புகளின் சார்புகளால் தரப்படுவதால், கூற்றுத் தருக்கவியல்
\emph{மெய்மதிப்புசார்} என்று கூறுகிறோம். குறிப்பாக, ஒரு கூற்று வெறுமனே
சூழ்நிலைசார் மெய்யாக இல்லாமல் இன்றியமையாமல் மெய்யா, அது மெய்யென்று
அறியப்பட்டுள்ளதா, அல்லது அது முன்பு மெய்யாக இருந்ததா பின்னர் மெய்யாகுமா
என்பதற்குப் பதிலாக இப்போது மெய்யா போன்ற மெய்மை--பொய்மை பற்றிய மேலதிக
வேறுபாடுகளைக் கூற்றுத் தருக்கவியலில் கருத்தில் கொள்ளமாட்டோம். மெய்
($\True$), பொய் ($\False$) என்ற இரு மெய்மதிப்புகளை மட்டுமே கருதுகிறோம்;
எனவே ஒரு கூற்று மெய்யுமல்ல பொய்யுமல்ல என்றோ அரைமெய் என்றோ அமையும்
சாத்தியத்தை விவாதத்திலிருந்து விலக்குகிறோம். மேலும், ஒரு !!a{formula}வின்
மெய்மதிப்பு அதன் பகுதிகளின் மெய்மதிப்புகளால் முற்றிலும் தீர்மானிக்கப்படும்
இணைப்புகளையே கவனிக்கிறோம்; அதன் பொருளால் தீர்மானிக்கப்படுவதை அல்ல.
குறிப்பாக, ஆங்கில நிபந்தனைக்கூற்றுகளின் மெய்மதிப்பு இவ்வகையில்
மெய்மதிப்புசார் தானா என்பது சர்ச்சைக்குரியது. பொருட்சார் நிபந்தனையான
$\lif$ மெய்மதிப்புசார்; மெய்மதிப்புசார் அல்லாத நிபந்தனைகளை வேறு
தருக்கவியல்கள் கையாளுகின்றன.

% SEGMENT OLP-0057-S02
மெய்மதிப்புசார் கூற்றுத் தருக்கவியலின் கோட்பாட்டையும் மீக்கோட்பாட்டையும்
வளர்க்க, முதலில் அதன் கோவைகளின் தொடரமைப்பையும் பொருண்மையையும் வரையறுக்க
வேண்டும். இணைப்புகளைப் பயன்படுத்தி !!{propositional variable}s இலிருந்து
!!{formula}s கட்டமைக்கும் ஒரு முறையை விவரிப்போம். வேறு வரையறைகளும்
சாத்தியம். பிற முறைமைகள் வேறு குறியீடுகளைத் தேர்ந்தெடுக்கலாம்; வேறு
இணைப்புக் கணங்களை மூல இணைப்புகளாகக் கொள்ளலாம்; அடைப்புக்குறிகளை வேறுபடப்
பயன்படுத்தலாம். போலந்து குறிமுறை எனப்படும் முறையில் அடைப்புக்குறிகளே
இல்லாமலும் இருக்கலாம். ஆனாலும் எல்லா அணுகுமுறைகளிலும், அமைப்பு விதிகள்
!!{formula}s இன் கணத்தை \emph{தொகுத்தறிதலாக} வரையறுக்கின்றன. அவை சரியாக
அமைக்கப்பட்டால், ஒவ்வொரு கோவையும் அமைப்பு விதிகளிலிருந்து அடிப்படையில்
ஒரே ஒரு வழியில் மட்டுமே உருவாகும். கோவைகளுக்குத் \emph{தனித்த
வாசிப்புத்தன்மையை} அளிக்கும் இந்தத் தொகுத்தறிதல் வரையறையால், அதே
முறையான தொகுத்தறிதல் வரையறையைப் பயன்படுத்தி அவற்றிற்குப் பொருள் தரலாம்.

% SEGMENT OLP-0057-S03
கோவைகளுக்குப் பொருள் தருவது பொருண்மையின் துறை. கூற்றுத் தருக்கவியலின்
பொருண்மையில் மையக் கருத்து, ஒரு !!a{valuation}வில் நிறைவுறுத்துதல்.
!!^a{valuation}~$\pAssign{v}$, !!{propositional variable}s க்கு
$\True$, $\False$ என்ற மெய்மதிப்புகளை ஒதுக்குகிறது. எந்த
!!a{valuation}வும் ஒவ்வொரு !!a{formula}~$!A$ க்கும் ஒரு மெய்மதிப்பு
$\pValue{v}(!A)$ ஐத் தீர்மானிக்கிறது. $\pValue{v}(!A) = \True$ எனில்,
அப்போதும் அப்போது மட்டுமே, !!^a{formula} ஆனது
!!a{valuation}~$\pAssign{v}$ இல் நிறைவுறுத்தப்படுகிறது; இதை
$\pSat{v}{!A}$ என்று எழுதுகிறோம். தருக்க இணைப்புகளுக்கான மெய்மதிப்புச்
சார்புகளைப் பயன்படுத்தி, எடுத்துக்காட்டாக $!A \land !B$ நிறைவுறுவது
$!A$, $!B$ ஆகியவை நிறைவுறுவது அல்லது நிறைவுறாதது ஆகியவற்றின் வழியாக,
இந்தத் தொடர்பையும் $!A$ இன் கட்டமைப்பின்மீதான தொகுத்தறிதலால் வரையறுக்கலாம்.

% SEGMENT OLP-0057-S04
வாக்கியங்களுக்கான நிறைவுறுத்தல் தொடர்பான $\pSat{v}{!A}$ ஐ அடிப்படையாகக்
கொண்டு, மெய்மம், பொருண்மைப் பின்விளைவு, நிறைவுறத்தக்க தன்மை ஆகிய அடிப்படை
பொருண்மைக் கருத்துகளை வரையறுக்கலாம். ஒவ்வொரு !!a{valuation}வும் ஒரு
!!a{formula}வை நிறைவுறுத்தினால், அதாவது எந்த~$\pAssign{v}$ க்கும்
$\pValue{v}(!A) = \True$ எனில், அந்த !!^a{formula} ஒரு மெய்மம்;
$\Entails !A$ என்று எழுதுகிறோம். $\Gamma$ இல் உள்ள எல்லா !!{formula}s
ஐயும் நிறைவுறுத்தும் ஒவ்வொரு !!a{valuation}வும் $!A$ ஐ நிறைவுறுத்தினால்,
!!{formula}s கொண்ட அந்தக் கணத்திலிருந்து அந்த வாய்பாடு பொருண்மைப்படிப்
பின்விளைகிறது;
$\Gamma \Entails !A$ என்று எழுதுகிறோம். ஏதோ ஒரு !!a{valuation} ஒரே
நேரத்தில் ஒரு !!{formula}s கணத்தின் எல்லா உறுப்புகளையும் நிறைவுறுத்தினால்,
அந்தக் கணம் நிறைவுறத்தக்கது. !!{formula}s தொகுத்தறிதலாக
வரையறுக்கப்படுவதாலும், !!{formula}s இன் கட்டமைப்பின்மீதான தொகுத்தறிதலால்
நிறைவுறுத்தல் வரையறுக்கப்படுவதாலும், நமது பொருண்மையின் பண்புகளை நிறுவவும்
வரையறுக்கப்பட்ட பொருண்மைக் கருத்துகளைத் தொடர்புபடுத்தவும் தொகுத்தறிதலைப்
பயன்படுத்தலாம்.


\end{document}
